% Esqueleto para redação do texto (mesmo formato SBPO 2026 ST).
% O PDF de submissão (sbpo2026.tex) deve ter no máximo 12 páginas; aqui não ativamos o aviso automático.
% Preencha \texttt{sections/esqueleto\_corpo.tex} ou copie blocos para \texttt{sbpo2026.tex}.
\documentclass[11pt,a4paper]{article}
% Pacotes e ajustes comuns (SBPO 2026 ST — margens, Times 11, PT-BR).
\usepackage[a4paper,top=3.3cm,bottom=2.5cm,left=2.9cm,right=2.9cm]{geometry}
\usepackage{mathptmx}
\usepackage[T1]{fontenc}
\usepackage[utf8]{inputenc}
\usepackage[brazil]{babel}
\usepackage{microtype}
\usepackage{amsmath,amssymb}
\usepackage{array}
\usepackage{booktabs}
\usepackage{graphicx}
\usepackage{pgfplots}
\pgfplotsset{compat=1.18}
\usepackage{tikz}
\usetikzlibrary{arrows.meta,positioning,calc}
\usepackage{hyperref}
\usepackage{enumitem}
\setlist{nosep}

\linespread{1}
\renewcommand{\normalsize}{\fontsize{11}{13.2}\selectfont}
\normalsize

\hypersetup{colorlinks=true,linkcolor=black,citecolor=black,urlcolor=blue}

% Limite SBPO 2026 (ST): 12 páginas no PDF final (texto + figuras + referências).
\newcount\sbpoMaxPages
\sbpoMaxPages=12
\newif\ifsbpoLimitPages
\sbpoLimitPagesfalse
\AtEndDocument{%
  \ifsbpoLimitPages
    \ifnum\value{page}>\sbpoMaxPages\relax
      \PackageWarning{sbpo2026}{%
        Limite de paginas SBPO ST: este PDF tem \arabic{page} pagina(s) (maximo \the\sbpoMaxPages).%
        \MessageBreak Reduza texto, equacoes, figuras/tabelas ou bibliografia; recompile e verifique o log.%
      }%
    \fi
  \fi
}

\sbpoLimitPagesfalse

\title{\textbf{[Título provisório]}\\[0.35em]\large Esqueleto de texto --- SBPO 2026 (ST)}
\date{}

\begin{document}
\maketitle

\begin{abstract}
\textit{[Resumo: contexto em 1--2 frases; gap; proposta/método em 1 frase; resultado principal em 1 frase; contribuição. Alvo típico: 150--250 palavras, conforme edital.]}
\end{abstract}

\noindent\textbf{Palavras-chave:}
\textit{[3--5 termos do edital / IEEE / área.]}

\tableofcontents
\newpage

% ---------------------------------------------------------------------------
% Skeleton body: replace \textit{[...]} blocks with final paragraphs.
% Labels match the main article (\texttt{sbpo2026.tex}) for easier merging.
% ---------------------------------------------------------------------------

\section{Introduction}
\label{sec:introducao}

\subsection{Context and problem}
\textit{[Why agile flow + human factors matter for OR, teaching, or industry. Opening citations: Kanban, Scrum, simulation.]}

\subsection{Objective}
\textit{[One clear sentence + 2--3 measurable or modeling sub-goals.]}

\subsection{Contributions}
\textit{[Numbered list or itemize: model; implementation; experiments; reproducibility. Avoid vague ``original contribution''.]}

\subsection{Paper organization}
\textit{[Standard paragraph: Sec.~\ref{sec:relacionados}--\ref{sec:conclusao} map to sections.]}

% ---------------------------------------------------------------------------
\section{Related work}
\label{sec:relacionados}

\subsection{Flow, Kanban, and software metrics}
\textit{[Poppendieck, Anderson, flow metrics, CFD if relevant.]}

\subsection{Simulation and serious games in software engineering}
\textit{[Educational simulators; DES; ABM if it fits; gap vs.\ this proposal.]}

\subsection{Synthesis and positioning}
\textit{[Optional table ``work | focus | limitation''; where this paper fits.]}

% ---------------------------------------------------------------------------
\section{System description and operational model}
\label{sec:modelo}

\subsection{Overview and entities}
\textit{[Members, cards, columns, flow order. Optional board figure.]}

\subsection{Global parameters}
\textit{[Sprint, WIP, PRNG, coefficients $\beta$, $\gamma$, etc.; pointer to parameter table when available.]}

\subsection{Synergy and traits}
\textit{[Notação $s_{ij}$; sinergia par a par; ver texto final do artigo.]}

\subsection{Calendar and time advance}
\textit{[Global day, sprint, ceremony; what changes each step.]}

\subsection{Planning and pull policies}
\textit{[Planning $\rightarrow$ analysis; backlog pull and WIP / $P_{\max}$ limits.]}

\subsection{Daily capacity and column work}
\textit{[Dice draw, specialist, roles; capacity spend per column.]}

\subsection{Collaboration in \texttt{dev} and handoffs}
\textit{[Multipliers, clamp, rework; main equations without over-implementing.]}

\subsection{Pseudocode or flowchart (optional)}
\textit{[Pointer to Figure~\ref{fig:pseudo} when included.]}

% ---------------------------------------------------------------------------
\section{Scrum ceremonies on the calendar}
\label{sec:ritos}

\subsection{Ceremony $\leftrightarrow$ day mapping}
\textit{[Table~\ref{tab:ritos}: planning, daily, review, retro.]}

\subsection{Effects on the engine}
\textit{[What each ceremony changes in state; random events if any.]}

% ---------------------------------------------------------------------------
\section{Prototype architecture and reproducibility}
\label{sec:arquitetura}

\subsection{Modules and responsibilities}
\textit{[Engine, UI, log export; language and repository.]}

\subsection{Text--code traceability}
\textit{[How type/function names appear in the PDF; seed and result JSON.]}

\subsection{Flow figure (optional)}
\textit{[Figure~\ref{fig:fluxoQuadro}.]}

% ---------------------------------------------------------------------------
\section{Experimental design}
\label{sec:experimentos}

\subsection{Objective and hypotheses}
\textit{[Operational H1, H2, H3; what is fixed vs.\ varied.]}

\subsection{Team, backlog, and scenarios A/B/C/D}
\textit{[Short description; pointer to Table~\ref{tab:param}.]}

\subsection{Metrics and replication protocol}
\textit{[CFD, cycle time, throughput; command to regenerate figures.]}

% ---------------------------------------------------------------------------
\section{Results}
\label{sec:resultados}

\subsection{Time series and CFD}
\textit{[Figure~\ref{fig:cfdA} or per-scenario variants.]}

\subsection{Cycle time and throughput}
\textit{[Figure~\ref{fig:bar}, Figure~\ref{fig:tp}; Table~\ref{tab:resumoCenarios}.]}

\subsection{Qualitative reading}
\textit{[2--4 sentences linking numbers to hypotheses; do not over-claim beyond the design.]}

% ---------------------------------------------------------------------------
\section{Discussion}
\label{sec:discussao}

\subsection{Model and experiment limitations}
\textit{[Known simplifications; threats to external validity.]}

\subsection{Future work}
\textit{[Calibration on real data; more traits; alternative policies.]}

% ---------------------------------------------------------------------------
\section{Conclusion}
\label{sec:conclusao}

\textit{[Restate problem + method + main finding + contribution; one paragraph plus optional closing line.]}


% Esqueleto: sem bibliografia obrigatória; descomente ao citar.
% \bibliographystyle{plain}
% \bibliography{bibliografia}

\end{document}
